\chapter{Sistemas Relacionados}
\label{cap:sistemas}

En este capítulo analizo el panorama actual de sistemas basados en modelos de lenguaje grandes (LLMs) orientados a la gestión empresarial. Estudio tanto las soluciones comerciales de las grandes tecnológicas como las alternativas de código abierto que persiguen objetivos similares a los de mi proyecto. Esto me permite situar el sistema que he desarrollado en su contexto adecuado y justificar las decisiones de diseño que he adoptado.

% ============================================================
\section{Asistentes Comerciales Basados en LLM}
% ============================================================

\subsection{ChatGPT y GPT Enterprise (OpenAI)}

ChatGPT es el asistente conversacional más conocido del mercado, desarrollado por OpenAI. Su versión empresarial, ChatGPT Enterprise, incorpora características orientadas al ámbito corporativo: cifrado en tránsito y en reposo, garantía de no utilización de datos para entrenamiento del modelo, panel de administración para gestión de usuarios y cumplimiento SOC~2.

Sin embargo, ChatGPT Enterprise sigue siendo un servicio en la nube: todos los datos ---consultas, documentos subidos y respuestas--- se procesan en servidores de OpenAI. Aunque la compañía asegura que los datos empresariales no se utilizan para re-entrenamiento, el mero hecho de que la información salga del perímetro corporativo supone un riesgo inaceptable en muchos sectores regulados. Adicionalmente, el modelo de precios por token procesado puede escalar rápidamente en organizaciones con un uso intensivo.

\subsection{Microsoft Copilot para Microsoft 365}

Microsoft Copilot integra modelos GPT-4 de OpenAI directamente en las aplicaciones de Microsoft 365 (Word, Excel, Teams, SharePoint). Su ventaja principal es la integración nativa con el ecosistema Microsoft: accede a los documentos de SharePoint, los correos electrónicos de Outlook y los mensajes de Teams del usuario para contextualizar las respuestas.

No obstante, Copilot presenta varias limitaciones para un despliegue corporativo con requisitos de soberanía:

\begin{itemize}
    \item \textbf{Dependencia de la nube de Microsoft:} Todo el procesamiento se realiza en Azure, sin opción de ejecución on-premise.
    \item \textbf{Coste por usuario:} El precio de Copilot para Microsoft 365 es de 30 USD/usuario/mes, lo que resulta prohibitivo para organizaciones grandes.
    \item \textbf{Caja negra:} No es posible modificar la lógica de enrutamiento, personalizar el comportamiento del agente ni integrar fuentes de información externas no contempladas por Microsoft.
    \item \textbf{Sin control sobre el modelo:} No se puede elegir el modelo, ajustar parámetros de inferencia ni realizar fine-tuning.
\end{itemize}

\subsection{Google Gemini y Vertex AI}

Google ofrece sus modelos Gemini tanto como producto de consumo (Gemini en Google Workspace) como plataforma de desarrollo empresarial (Vertex AI). Vertex AI permite desplegar modelos propios de Google junto con modelos de terceros, con opciones de RAG integradas a través de su servicio Vertex AI Search.

Las mismas limitaciones fundamentales aplican: procesamiento en la nube de Google, dependencia de un proveedor, costes recurrentes por uso y ausencia de control total sobre los datos. Aunque Google ofrece opciones de localización de datos en regiones específicas, no soporta despliegue on-premise.

\subsection{Amazon Bedrock}

Amazon Bedrock es la plataforma de AWS que ofrece acceso a modelos de lenguaje de múltiples proveedores (Anthropic Claude, Meta LLaMA, Mistral, Amazon Titan) mediante una API unificada. Incluye capacidades de RAG a través de Knowledge Bases y agentes con herramientas.

Bedrock ofrece mayor flexibilidad que los productos de Microsoft o Google en cuanto a elección de modelo, pero sigue siendo un servicio en la nube con las mismas implicaciones de soberanía de datos y costes por uso.

% ============================================================
\section{Soluciones de Código Abierto para RAG}
% ============================================================

\subsection{PrivateGPT}

PrivateGPT\footnote{\url{https://github.com/zylon-ai/private-gpt}} es un proyecto de código abierto que permite interactuar con documentos propios utilizando LLMs locales. Su enfoque es similar al de este proyecto en la filosofía de privacidad: todo se ejecuta en local sin enviar datos a servicios externos.

Sin embargo, PrivateGPT se posiciona como una herramienta individual más que como un sistema empresarial: carece de gestión de usuarios, no ofrece autenticación corporativa (SSO), no incluye integración con SharePoint ni con fuentes legislativas, no dispone de monitorización de GPU y su capacidad de extensión es limitada.

\subsection{LocalAI}

LocalAI\footnote{\url{https://github.com/mudler/LocalAI}} es un servidor de API compatible con OpenAI que permite ejecutar modelos localmente. A diferencia de Ollama, LocalAI soporta una gama más amplia de backends (llama.cpp, GPT4All, Diffusers) pero con mayor complejidad de configuración.

LocalAI se centra exclusivamente en la capa de servicio del modelo (equivalente a Ollama + LiteLLM); no incluye pipeline RAG, interfaz de usuario, sistema de ingesta ni agente enrutador, por lo que solo cubriría una fracción de las funcionalidades del sistema desarrollado.

\subsection{AnythingLLM}

AnythingLLM\footnote{\url{https://github.com/Mintplex-Labs/anything-llm}} es una solución de código abierto que combina interfaz de chat, capacidad RAG y gestión de documentos en una aplicación unificada. Soporta múltiples proveedores de LLM (OpenAI, Ollama, LM Studio) y varias bases de datos vectoriales.

Si bien AnythingLLM se acerca más a la visión integral del sistema desarrollado en este proyecto, presenta diferencias importantes: no dispone de agente enrutador inteligente (toda consulta se procesa de la misma forma), no incluye integración con SharePoint, no ofrece OCR para documentos escaneados, no integra búsqueda en internet ni consulta legislativa del BOE, y no incorpora monitorización de infraestructura.

\subsection{Danswer / Onyx}

Danswer\footnote{\url{https://github.com/danswer-ai/danswer}} (recientemente renombrado a Onyx) es un sistema RAG empresarial de código abierto que se conecta a múltiples fuentes de información (Confluence, Slack, Google Drive, SharePoint). Incluye interfaz web, gestión de usuarios y capacidades de búsqueda.

Danswer es posiblemente el proyecto de código abierto más comparable al desarrollado en este TFG. Sin embargo, su arquitectura depende de APIs externas para la generación (por defecto OpenAI), no incluye ejecución local de modelos, y su personalización requiere modificar el código fuente del proyecto en lugar de utilizar un mecanismo de extensión limpio como los Pipelines de OpenWebUI.

\subsection{Frameworks: LangChain y LlamaIndex}

LangChain\footnote{\url{https://langchain.com/}} \cite{langchainRetrievalDocs} y LlamaIndex\footnote{\url{https://www.llamaindex.ai/}} \cite{llamaindexRagDocs} son frameworks de desarrollo para construir aplicaciones con LLMs. No son productos finales, sino herramientas de programación que proporcionan abstracciones para tareas como retrieval, chunking, cadenas de procesamiento y agentes.

Estos frameworks son útiles como componentes dentro de un sistema mayor, pero por sí solos no constituyen una solución: no proporcionan interfaz de usuario, autenticación, ingesta automatizada, monitorización ni despliegue. En este proyecto, se utilizan selectivamente conceptos y patrones de LangChain sin adoptar el framework como dependencia completa, priorizando implementaciones directas con menor acoplamiento.

% ============================================================
\section{Análisis Comparativo}
% ============================================================

La Tabla~\ref{tab:comparativa-sistemas} presenta una comparación exhaustiva entre el sistema que he desarrollado y las alternativas actuales, evaluando las funcionalidades clave para un despliegue corporativo con requisitos de privacidad y autonomía.

\begin{table}[H]
\centering
\caption{Comparación del sistema JARVIS con soluciones alternativas.}
\label{tab:comparativa-sistemas}
\resizebox{\textwidth}{!}{%
\begin{tabular}{lccccccc}
\toprule
\textbf{Característica} & \textbf{JARVIS} & \textbf{ChatGPT Ent.} & \textbf{Copilot 365} & \textbf{PrivateGPT} & \textbf{AnythingLLM} & \textbf{Danswer} & \textbf{LocalAI}\\
\midrule
On-Premise 100\%       & \checkmark & $\times$   & $\times$    & \checkmark & \checkmark & Parcial    & \checkmark \\
Código abierto         & \checkmark & $\times$   & $\times$    & \checkmark & \checkmark & \checkmark & \checkmark \\
UI conversacional      & \checkmark & \checkmark & \checkmark  & \checkmark & \checkmark & \checkmark & $\times$ \\
Multi-fuente RAG       & \checkmark & Parcial    & Parcial     & Parcial    & Parcial    & \checkmark & $\times$ \\
SharePoint             & \checkmark & $\times$   & \checkmark  & $\times$   & $\times$   & \checkmark & $\times$ \\
OCR integrado (GPU)    & \checkmark & $\times$   & $\times$    & $\times$   & $\times$   & $\times$   & $\times$ \\
Web Scraping           & \checkmark & $\times$   & Parcial     & $\times$   & $\times$   & $\times$   & $\times$ \\
Búsqueda en Internet   & \checkmark & \checkmark & \checkmark  & $\times$   & $\times$   & $\times$   & $\times$ \\
Agente enrutador       & \checkmark & Interno    & Interno     & $\times$   & $\times$   & $\times$   & $\times$ \\
SSO corporativo        & \checkmark & \checkmark & \checkmark  & $\times$   & $\times$   & \checkmark & $\times$ \\
Consulta legislativa   & \checkmark & $\times$   & $\times$    & $\times$   & $\times$   & $\times$   & $\times$ \\
Monitorización GPU     & \checkmark & $\times$   & $\times$    & $\times$   & $\times$   & $\times$   & $\times$ \\
Fine-tuning local      & \checkmark & $\times$   & $\times$    & $\times$   & $\times$   & $\times$   & Parcial \\
Extensible (plugins)   & \checkmark & $\times$   & $\times$    & $\times$   & $\times$   & Parcial    & $\times$ \\
Coste por token        & 0 \euro{}  & \$\$\$     & 30 \$/usr   & 0 \euro{}  & 0 \euro{}  & Variable   & 0 \euro{} \\
\bottomrule
\end{tabular}%
}
\end{table}

% ============================================================
\section{Posicionamiento del Sistema Desarrollado}
\label{sec:posicionamiento}
% ============================================================

Tras este análisis comparativo, concluyo que el sistema JARVIS que he construido ocupa un nicho diferenciado que no está cubierto por ninguna de las alternativas analizadas de forma individual:

\begin{enumerate}
    \item \textbf{Frente a soluciones comerciales} (ChatGPT Enterprise, Copilot, Vertex AI, Bedrock): el sistema ofrece la misma experiencia conversacional pero con soberanía total sobre los datos. Ningún documento, consulta o respuesta abandona el perímetro del servidor corporativo. El coste operativo se limita a la electricidad y el hardware, sin costes recurrentes por token procesado.

    \item \textbf{Frente a soluciones de código abierto} (PrivateGPT, LocalAI, AnythingLLM): el sistema desarrollado ofrece un nivel de integración y funcionalidad significativamente superior. Mientras que estas alternativas se centran típicamente en una única capacidad (chat local, RAG básico), JARVIS integra múltiples fuentes de datos (SharePoint, web, BOE), múltiples modos de operación (RAG, scraping, OCR, búsqueda web, visión), autenticación corporativa, monitorización de infraestructura y un agente enrutador inteligente.

    \item \textbf{Frente a Danswer/Onyx}: es la alternativa más comparable en ambición. Sin embargo, Danswer depende de APIs externas para la generación, mientras que el sistema desarrollado ejecuta todo localmente. Además, la extensibilidad mediante Pipelines de OpenWebUI ofrece un mecanismo más limpio de personalización que modificar directamente el código fuente de Danswer.

    \item \textbf{Frente a frameworks} (LangChain, LlamaIndex): estos no son productos finales sino herramientas de desarrollo. El sistema JARVIS es un sistema desplegable llave en mano, con interfaz de usuario, autenticación, monitorización y documentación de despliegue incluidas de serie.
\end{enumerate}

En resumen, la contribución principal de mi trabajo respecto al estado actual es ofrecer un \textbf{sistema RAG corporativo completo, on-premise, extensible y con múltiples fuentes de datos}, combinando la calidad de la experiencia de usuario de los productos comerciales con la privacidad y autonomía del código abierto.
